\section{Notes on Uninformed and Informed Bandit Feedback}\label{h1a:uninformed-informed}

Existing literature on bandit-feedback learning in games exhibits a subtle but critical difference in the feedback model. For instance, the self-play framework in \citet{ito2025instancedependent} assumes the learner observes only the reward of the realized actions.
In contrast, the algorithms proposed by \citet{odonoghue2021matrix}, including variants of UCB, Thompson Sampling, and K-learning, operate under the assumption that the learner observes the reward and, in addition, the action taken by the opponent.
As this distinction is rarely addressed by prior work,
we introduce the terms ``uninformed'' and ``informed'' to formally differentiate between these two settings based on the availability of the opponent's action.

A similar distinction on informedness exists in the broader field of learning Markov games.
A Markov game is a game where the utility depends on the action of players and a state, and the state is changed every round in a Markovian way determined by the players' actions.
Our setting is a special case of Markov games with only one state.
Post-hoc observations on opponents' actions is assumed by some research \citep{wei2017online, bai2020provable, xie2020learning, liu2021sharp},
but in some other work, this information is hidden \citep{tian2021online, mao2022improving, cai2023uncoupled, appel2025regret}.
Out of these papers, \citet{cai2023uncoupled} define ``uncoupled'' learning, and \citet{mao2022improving} use the term ``decentralized'' learning, both of which imply the lack of observation of the opponent's actions (our ``uninformed'' feedback).


